\documentclass[11pt, letterpaper]{article}
\usepackage{mobi-lectura}

\title{Problema del cumpleaños: probabilidad de coincidencia en un grupo de $n$ personas}
\author{Alejandra Tabares Pozos}
\date{}

\begin{document}

\maketitle
\tableofcontents
\newpage

\section{Planteamiento y modelo (generalización a $d$ días equiprobables)}

Sea $d\in\mathbb{N}$ el número de días posibles (por ejemplo, $d=365$ si ignoramos años bisiestos).
Modelamos el cumpleaños de cada persona como una variable aleatoria discreta
\[
X_i \in \{1,2,\dots,d\}, \qquad i=1,\dots,n,
\]
bajo los supuestos estándar:
\begin{itemize}
\item \textbf{Equiprobabilidad:} $\mathbb{P}(X_i=j)=\frac{1}{d}$ para todo $j\in\{1,\dots,d\}$.
\item \textbf{Independencia:} $X_1,\dots,X_n$ son independientes (i.i.d.).
\end{itemize}

El espacio muestral es $\Omega=\{1,\dots,d\}^n$, con distribución uniforme
\[
\mathbb{P}((x_1,\dots,x_n))=\frac{1}{d^n}.
\]

El evento de interés es
\[
A=\{\exists\, i\neq j:\ X_i=X_j\},
\]
es decir, \emph{al menos dos personas comparten la misma fecha de nacimiento}.
La estrategia central es trabajar con el complemento
\[
A^c=\{X_1,\dots,X_n\ \text{son todos distintos}\}.
\]

\subsection*{Interpretación secuencial (clave para el conteo)}
Piense en asignar cumpleaños de forma secuencial:
\begin{itemize}
\item La persona 1 puede caer en cualquiera de los $d$ días.
\item Para que no haya repetición, la persona 2 debe caer en uno de los $d-1$ días restantes.
\item La persona 3, en uno de los $d-2$ días restantes.
\item En general, la persona $k$ (contando desde $k=1$) debe caer en uno de los $d-(k-1)$ días restantes.
\end{itemize}
Este razonamiento convierte un evento global (``todos distintos'') en una multiplicación de probabilidades condicionales.

\section{Cálculo exacto por probabilidad complementaria (caso general $d$)}

\subsection{Probabilidad de que todos sean distintos}
Por la regla del producto:
\begin{align}
\mathbb{P}(A^c)
&= \mathbb{P}(\text{persona 2 distinta de persona 1})\cdot
\mathbb{P}(\text{persona 3 distinta de las anteriores}\mid \dots)\cdots \nonumber\\
&= \frac{d-1}{d}\cdot \frac{d-2}{d}\cdots \frac{d-(n-1)}{d}
= \prod_{k=0}^{n-1}\left(1-\frac{k}{d}\right).
\end{align}
En notación de factorial descendente, $(d)_n=d(d-1)\cdots(d-n+1)$,
\[
\mathbb{P}(A^c)=\frac{(d)_n}{d^n}.
\]
Equivalentemente, para $n\le d$,
\[
(d)_n=\frac{d!}{(d-n)!}
\quad\Longrightarrow\quad
\mathbb{P}(A^c)=\frac{d!}{(d-n)!\,d^n}.
\]

\subsection{Probabilidad de al menos una coincidencia}
Finalmente,
\begin{equation}
\boxed{
\mathbb{P}(A)=1-\mathbb{P}(A^c)
=1-\prod_{k=0}^{n-1}\left(1-\frac{k}{d}\right)
=1-\frac{(d)_n}{d^n}.
}
\end{equation}
Si $n>d$, por el principio del palomar, necesariamente hay repetición y
\[
\mathbb{P}(A)=1.
\]

\section{Conteo explícito: por qué el problema ``explota''}

La razón combinatoria detrás de las fórmulas anteriores es:
\begin{itemize}
\item Total de asignaciones posibles de cumpleaños (con repetición): $|\Omega| = d^n$.
\item Asignaciones \emph{sin repetición} (todos distintos): primero elegir $n$ días distintos y asignarlos a $n$ personas,
lo cual produce $d(d-1)\cdots(d-n+1)=(d)_n$ posibilidades.
\item Asignaciones con \emph{al menos una repetición}: $d^n-(d)_n$.
\end{itemize}
Por tanto,
\[
\mathbb{P}(A)=\frac{d^n-(d)_n}{d^n}.
\]

\subsection{Ejemplo pequeño que muestra el conteo (y su diferencia)}
Tome $d=5$ días y $n=3$ personas.
\[
d^n = 5^3 = 125.
\]
Las asignaciones sin repetición son
\[
(d)_n = 5\cdot 4\cdot 3 = 60.
\]
Entonces las asignaciones con al menos una repetición son
\[
d^n-(d)_n = 125-60 = 65,
\]
y por ende
\[
\mathbb{P}(A)=\frac{65}{125}=0.52,
\qquad
\mathbb{P}(A^c)=\frac{60}{125}=0.48.
\]
Este ejemplo deja ver el fenómeno: \emph{aunque $n$ es pequeño, el número de pares potenciales crece rápidamente y las colisiones aparecen pronto.}

\subsection{Escala real: $d=365$ y $n=23$}
Con $d=365$ y $n=23$:
\[
d^n = 365^{23} \approx 8.565\times 10^{58}
\quad (\text{un entero de }59\text{ dígitos}).
\]
Mientras que
\[
(d)_n = 365\cdot 364\cdot \ldots \cdot 343 \approx 4.220\times 10^{58}
\quad (\text{también }59\text{ dígitos}).
\]
La diferencia
\[
d^n-(d)_n \approx 4.345\times 10^{58}
\]
es del mismo orden de magnitud que los términos involucrados: esto ilustra por qué
\emph{un conteo directo ``listando casos'' es imposible}, y por qué la fórmula del producto es la vía correcta.

Numéricamente (usando la fórmula exacta),
\[
\mathbb{P}(A)=1-\prod_{k=0}^{22}\left(1-\frac{k}{365}\right)\approx 0.5073.
\]
Es decir, con $23$ personas ya es ligeramente mayor al $50\%$.

\section{Aproximación útil para $d$ grande (opcional pero muy informativa)}
Partimos de
\[
\ln\mathbb{P}(A^c)=\sum_{k=0}^{n-1}\ln\left(1-\frac{k}{d}\right).
\]
Para $k/d$ pequeño, $\ln(1-x)\approx -x$, entonces
\[
\ln\mathbb{P}(A^c)\approx -\sum_{k=0}^{n-1}\frac{k}{d}
= -\frac{n(n-1)}{2d}.
\]
Así,
\[
\boxed{
\mathbb{P}(A)\approx 1-\exp\!\left(-\frac{n(n-1)}{2d}\right).
}
\]
Esta aproximación sugiere el umbral de ``sorpresa'': cuando $\frac{n(n-1)}{2d}\approx 1$,
la probabilidad de colisión deja de ser pequeña.

\section*{Referencias}
\begin{enumerate}[label={[\arabic*]}]
    \item Feller, W. (1968). \textit{An Introduction to Probability Theory and Its Applications}, Vol.~1, 3rd ed. Wiley. Cap.~2.
    \item Ross, S.\,M. (2014). \textit{Introduction to Probability Models}, 11th ed. Academic Press. Cap.~1.
    \item DeGroot, M.\,H. \& Schervish, M.\,J. (2012). \textit{Probability and Statistics}, 4th ed. Pearson. Cap.~1.
\end{enumerate}

\end{document}
